%======================================================================
\section{Discussion and questions}\label{sec:discussion}

\paragraph{One presentation for every prime.} The algebra of Theorem~A and the groups of Theorem~B come from one
presentation, fixed before the prime. The only place where the prime enters the construction is detection. There,
central elements are separated in finite quotients using long-word phases modulo integers coprime to $p$. The phases
of the hierarchy live modulo products of the encoding lengths $L_j=2^{h_j+1}+1$, and these products may share factors
with $p$. Reducing the phases modulo the part prime to $p$, which is unbounded (Section~\ref{sec:detection}), removes
this dependence. For example, an odd prime divides some $2^m+1$ only if the multiplicative order of $2$ modulo $p$ is
even, so primes such as $7$, $23$, $31$ and $47$ never divide any encoding length.

\paragraph{Exponents.} The groups $G_p$ have unbounded exponent, since a finitely generated residually finite group of
exponent dividing $p^a$ is finite by Zel'manov's theorems~\cite{ZelmanovOdd,ZelmanovTwo}. Whether an infinite
finitely presented group can have finite exponent remains open, and residual finiteness rules out the present method
for that question.

\paragraph{Characteristic zero.} Theorem~C shows that presentations by unsigned word relations, which carry the
positive-characteristic construction, can never give a nil algebra that is not nilpotent in characteristic zero. One
might hope to replace the symmetric sum of Lemma~\ref{lem:symmetric} by an alternating one, using letters that
anticommute. This does not work as stated. Exchanging two whole chunks of lengths $r$ and $s$ then contributes the sign
$(-1)^{rs}$, so chunks of even length still commute and the orbit sums do not become alternating. A characteristic-zero
analogue would need a genuinely different cancellation identity, valid for chunks of arbitrary length, together with an
invariant that keeps arbitrarily long words nonzero.

\paragraph{Questions.}
\begin{enumerate}[label=(\arabic*),leftmargin=*]
\item Is there a finitely presented nil algebra over $\Q$ that is not nilpotent? By Theorem~C its relations must involve
  signs or coefficients.
\item The rank $12$ is inherited from the stability range of the argument and is not optimized. What is the least $n$
  for which $\St_n(R_p)$ has the properties of Theorem~B?
\end{enumerate}

%======================================================================
\section{How this work was done}\label{sec:methods}

\paragraph{Roles.} This paper is the work of its two authors together with AI models from more than one developer.
\begin{itemize}[leftmargin=*]
\item The authors set the direction throughout: they selected the problem from a review of OpenAI's release, fixed the
  scope (every prime, with characteristic zero treated separately), set the verification protocol, decided which lines to
  pursue and which to drop, and take full responsibility for every claim.
\item The models did most of the exploration, derivation, computation and drafting. In particular they found:
  \begin{itemize}
  \item a cancellation bound in characteristic $p$, later found to be a weaker form of Zalesskii's
    theorem~\cite{Zalesskii1985};
  \item the failure of the unordered-frame argument and its ordered-frame replacement;
  \item the sign-corrected group calculations;
  \item the treatment of phases for odd primes;
  \item Theorem~C.
  \end{itemize}
\end{itemize}

\paragraph{Verification.} The extension was worked out in three independent lanes, each with three rounds, and a separate
coordinating instance re-derived the repairs on which the argument rests.
\begin{itemize}[leftmargin=*]
\item \emph{The lanes.} One lane covered the hierarchy, cancellation and phases, and one covered the group-theoretic
  steps. The third was a refutation-first adversary, run by a model from a different developer, which tried to break
  every step and then audited the proofs line by line.
\item \emph{What verification found.} It established that the direct port of the unordered-frame stability argument
  fails over odd fields (Proposition~\ref{prop:unordered-fails}), which led to the ordered-frame replacement.
\item \emph{Exact computation.} Finite identities were checked exactly, in integer arithmetic modulo $p$:
  \begin{itemize}
  \item the symmetric-sum bound of Lemma~\ref{lem:symmetric}, and Zalesskii's degree $dp$, exhaustively on tuples of
    matrix units for small $p$ and $d$;
  \item the matrices of Proposition~\ref{prop:unordered-fails} and of the face comparison in
    Theorem~\ref{thm:stability};
  \item the coordinate changes of Section~\ref{sec:torsion};
  \item the valuation bound of Lemma~\ref{lem:det-pprime}.
  \end{itemize}
  The script, \texttt{verify/checks.py}, is part of the release.
\item \emph{A blind check.} The four load-bearing steps were then checked independently by a model from another developer
  that had not seen our proofs:
  \begin{itemize}
  \item the characteristic-$p$ cancellation and its propagation;
  \item the ordered-frame stability argument;
  \item the ordinary finite presentation with $p$-power root relations;
  \item the detection of central elements.
  \end{itemize}
  It accepted all four, and confirmed by exact computation that the characteristic-two bound of
  \cite[Lemma~3.2]{OAI247periodic} fails over $\F_3$, as Zalesskii's theorem predicts, so a bound depending on $p$ is needed. It
  independently flagged two characteristic-two steps in the imported inputs:
  \begin{itemize}
  \item the permutation lemma~\cite[Lemma~3.3]{OAI247rf};
  \item the swaps in the finite-rank inputs of~\cite[Section~5]{OAI247periodic}.
  \end{itemize}
  Both are replaced in Sections~\ref{sec:stability}--\ref{sec:detection}. It also observed that reducing the phases to
  moduli prime to $p$ removes the need for prime-dependent parameters, which we adopt. It judged Theorem~C correct and
  elementary.
\end{itemize}
The OpenAI manuscripts were read at a fixed commit of their public repository. OpenAI's Lean development formalizes the
characteristic-two periodic group of~\cite{OAI247periodic} and nothing in the present paper. No argument here has been
formally verified; the proofs are given in full, and the checks support them but do not replace them. The
inputs~\cite{OAI247periodic,OAI247rf} were produced by an OpenAI model, as OpenAI states, and have not been refereed.

\paragraph{Priority search.} On 15 October 2026 we searched arXiv, reading math.GR and math.RA newest first and running
full-text queries on finitely presented periodic groups, torsion groups and $p$-groups, nil algebras and Steinberg groups.
We also searched OpenAI's public repository and the general web. We found no infinite finitely presented periodic group,
and no infinite finitely presented residually finite $p$-group, for any odd prime $p$. The nearest recent work is
Ivanov's extension to every odd prime of another construction from the same release, the torsion-free group whose
group algebra over $\F_2$ has zero divisors~\cite{Ivanov2026}. It also replaces a characteristic-two cancellation by a
signed one, but it concerns group algebras and is independent of the present paper. For Theorem~C we found no explicit
statement in the literature on semigroup algebras, and, as said after its proof, we do not claim it as new.
